
\begin{remark}
Recall that we assume that $\mbg$ is a finite, weighted, directed
graph, without loops, sources or sinks (\cf
Assumption~\ref{digraph}). In particular, its Laplacian matrix $L$ is
a $\cb$ matrix and, since $\mbg$ is strongly connected, then $L$ is an
$\icb$ matrix (see Subsection~\ref{dictionary}). By
Proposition~\ref{chain-complex}, $\mcc_{\mathbb{G}}$ is a chain
complex of free $\mba$-modules. Moreover, we suppose in $\mba=\mbk[x]$
the $\mbn$-grading given by $\nu(L)$ (\cf Notation~\ref{vectornu} and
Assumption~\ref{grading}). Furthermore, we assume that $\mbg$ has a
$(\omega,\delta)$-enumeration or, equivalently, that $L$ is in
$\delta$-block echelon form (see \ref{echelon-omega}). Finally, we
suppose that $\mcc_0=\mbk[x]$ is endowed with the $\wrlo$
(Assumption~\ref{wrlo}).

Recall also that $\dif_0:\mcc_0=\mbk[x]\to \mbk[x]/I(\mcl)$ is defined
to be the natural projection onto the quotient ring (see
Definition~\ref{dif0}). Observe that we may interpret $\mcc_n=0$ and
$\dif_n=0$.
\end{remark}

Let us define a monomial ordering on each module $\mcc_{k}$ of the
$\cyc$ complex.

\begin{remark}\label{mon-ordering}
Recall that the elements of the enumerated free basis $\mcb_k$ of the
free module $\mcc_k$ are denoted by $e_{k,1},\ldots ,e_{k,r_k}$, where
$r_k=\rank(\mcc_k)=\lvert\cyc_{n,k+1}\rvert$. Their images under
$\dif_k$ are denoted by $f_{k-1,1},\ldots ,f_{k-1,r_k}$, with
$f_{k-1,j}:=\dif_k(e_{k,j})$ (Definition~\ref{cyccomplex}). Let
\begin{eqnarray*}
G_{k-1}:=\dif_k(\mcb_k)= \{\dif_{k}(e_{k,1}),\ldots
,\dif_{k}(e_{k,r_k})\}= \{f_{k-1,1},\ldots
,f_{k-1,r_k}\}\subset\mcc_{k-1}.
\end{eqnarray*}
By Assumption~\ref{wrlo}, we endow $\mcc_0=\mbk[x]$ with the
$\wrlo$. Using Definitions~\ref{secondsetofdefinitions}, where
$G_0=\dif_1(\mcb_1)$, with $\mcb_1$ a basis of $\mcc_1$ and
$\dif_1:\mcc_1\to\mcc_0$, we endow $\mcc_1$ with the monomial ordering
induced by the one in $\mcc_0$ and the subset $G_0$. Taking into
account that $G_{k-1}=\dif_k(\mcb_{k})$, with $\mcb_{k}$ a basis of
$\mcc_k$ and $\dif_k:\mcc_k\to\mcc_{k-1}$, we can proceed recursively,
and endow each $\mcc_k$ with the monomial ordering induced by the one
in $\mcc_{k-1}$ and the subset $G_{k-1}$.
\end{remark}

\begin{notation}\label{m-basis-elements}
With the notations as above in Remark~\ref{mon-ordering}, set
\begin{eqnarray*}
m^k_{j,i}:=\frac{\Lcm(\lm(\dif_{k}(e_{k,j})),\lm(\dif_k(e_{k,i})))}
{\lt(\dif_k(e_{k,i}))}=
\frac{\Lcm(\lm(f_{k-1,j}),\lm(f_{k-1,i}))}{\lt(f_{k-1,i})}.
\end{eqnarray*}
Similarly to Notation~\ref{module-quotients}, $M_i(G_{k-1})$ are the
monomial ideals of $\mba=\mbk[x]$ defined as follows:
$M_1(G_{k-1})=0$, and for each $i=2,\ldots,r_{k-1}$,
\begin{eqnarray*}
&& M_i(G_{k-1})=(\lt(f_{k-1,1}),\ldots,\lt(f_{k-1,i-1})):\lt(f_{k-1,i})=
  \\ && \langle\lt(f_{k-1,1})\rangle:\lt(f_{k-1,i})+\cdots+
  \langle\lt(f_{k-1,i-1})\rangle:\lt(f_{k-1,i})= (m^k_{1,i},\ldots
  ,m^k_{i-1,i}).
\end{eqnarray*}
Note that, if $\lm(f_{k-1,j})$ and $\lm(f_{k-1,i})$ are two monomials
of $\mcc_{k-1}$ which are multiple of different basis elements, then
$\langle\lt(f_{k-1,j})\rangle:\lt(f_{k-1,i})=0$. In such a case, we
understand $m_{j,i}^k=0$ (see Definition~\ref{firstsetofdefinitions})
\end{notation}
 
\begin{notation}\label{hyphk} 
We will prove Theorem~\ref{main} by induction on the degree of the
component of the complex. We specify now the statements to be shown.

\underline{Statement $\mbfh_0$}:
\begin{itemize}
\item[$(a)$] $G_{0}=\dif_{1}(\mcb_{1})$ is a Gr\"obner basis of
  $\ker(\dif_{0})\subset\mcc_{0}$.
\item[$(b)$] $\lt(\dif_{1}(C,\overline{C}))=
  (-1)^0x^{C\to\overline{C}}(C\cup\overline{C})$, for each
  $(C,\overline{C})\in\mcb_{1}$.
\end{itemize}

\underline{Statement $\mbfh_k$}, for $1\leq k\leq n-2$ (recall that
$n\geq 3$):
\begin{itemize}
\item[$(a)$] $G_{k}=\dif_{k+1}(\mcb_{k+1})$ is a Gr\"obner basis of
  $\ker(\dif_{k})\subset\mcc_{k}$ w.r.t. the monomial ordering defined
  as in Remark~\ref{mon-ordering}.
\item[$(b)$] $\lt(\dif_{k+1}(I_1,\ldots
  ,I_{k+2}))=(-1)^{k}x^{I_{k+1}\to I_{k+2}}(I_1,\ldots ,I_{k+1}\cup
  I_{k+2})$, for each basis element $(I_1,\ldots
  ,I_{k+2})\in\mcb_{k+1}$.
\end{itemize}

\underline{Statement $\mbfh_{n-1}$}:
$\dif_{n-1}:\mcc_{n-1}\to\dif_{n-2}$ is injective.
\end{notation}

\begin{remark}\label{degree0}
We already know that $\mbfh_0$ holds. Indeed, by
Corollary~\ref{GBI(L)}, $G_0=\dif_1(\mcb_1)$ is a Gr\"obner basis of
${\rm Im}(\dif_1)=I(\mcl)=\ker(\dif_0)$. Moreover, by Lemma~\ref{lm},
\begin{eqnarray*}
\lt(\dif_1(C,\overline{C}))=\lt(f_C)=m_C=x^{C\to\overline{C}}=
(-1)^0x^{C\to\overline{C}}(C\cup\overline{C}),
\end{eqnarray*}
where, according to Definition~\ref{cyccomplex}, $(C\cup
\overline{C})\equiv ([n])$ is identified with the unit element of
$\mcc_0=\mbk[x]$.
\end{remark}

\begin{remark}\label{exactnessn-1}
If $\mbfh_{n-2}$ holds, then $\mbfh_{n-1}$ holds. Indeed, for each
$i=2,\ldots ,r_{n-1}$, the basis element $e_{n-1,i}=(I_1,\ldots
,I_n)\in\mcb_{n-1}$ is a cyclically ordered partition of $[n]$ into
$n$ blocks, with $n\in I_n$. Thus, $(I_1,\ldots ,I_{n})
=(\{i_1\},\ldots ,\{i_{n-1}\},\{n\})$, where $\{i_1,\ldots
,i_{n-1},n\}=[n]$. Any different basis element
$e_{n-1,j}\in\mcb_{n-1}$, with $1\leq j<i$, will be of the same form:
$e_{n-1,j}=(\{j_1\},\ldots ,\{j_{n-1}\},\{n\})\in\mcb_{n-1}$. By part
$(b)$ of $\mbfh_{n-2}$,
\begin{eqnarray*}
&&\lt(f_{n-1,i})= \lt(\dif_{n-1}(e_{n-1,i}))=\\&&
  (-1)^{n-2}x^{I_{n-1}\to\{n\}}(\{i_1\},\ldots
  ,\{i_{n-2}\},\{i_{n-1},n\})
\end{eqnarray*}
and, similarly,
\begin{eqnarray*}
&&\lt(f_{n-1,j})=\lt(\dif_{n-1}(e_{n-1,j}))=\\&&
(-1)^{n-2}x^{J_{n-1}\to\{n\}}(\{j_1\},\ldots
,\{j_{n-2}\},\{j_{n-1},n\}).
\end{eqnarray*}
One sees that, since $e_{n-1,i}$ and $e_{n-1,j}$ are different, then
$\lt(f_{n-1,i})$ and $\lt(f_{n-1},j)$ are multiple of different basis
elements. Thus
$\langle\lt(f_{n-1,j})\rangle:\lt(f_{n-1,i})=0$ and
$M_i(G_{n-1})=0$, for all $i=1,\ldots ,r_{n-1}$ (see
Notation~\ref{m-basis-elements}). By Theorem~\ref{algorithm}, we
deduce that $\ker(\dif_{n-1})=0$.
\end{remark}

\begin{remark}\label{exactness1ton-2}
If $\mbfh_k$ holds, for every $k=0,\ldots ,n-1$, then
$\mcc_{\mathbb{G}}$ is an exact complex. Indeed, by
Remark~\ref{degree0} above, ${\rm Im}(\dif_1)=\ker(\dif_0)$. For
$1\leq k\leq n-2$, by part $(a)$ of $\mbfh_k$,
\begin{eqnarray*}
{\rm Im}(\dif_{k+1})=\langle
\dif_{k+1}(\mcb_{k+1})\rangle=\langle \mbox{Gr\"obner basis of
}\ker(\dif_k)\rangle =\ker(\dif_{k}).
\end{eqnarray*}
Finally, by definition, $\dif_{n}=0$, so ${\rm Im}(\dif_{n})=0$. On
the other hand, $\dif_{n-1}$ is injective by $\mbfh_{n-1}$.
\end{remark}

\begin{purpose}\label{reduction}
In the light of Remarks~\ref{degree0}, \ref{exactnessn-1} and
\ref{exactness1ton-2}, the proof of the main theorem is reduced to
show ``$\mbfh_{k-1}\Rightarrow \mbfh_{k}$, for $k=1,\ldots ,n-2$''.
\end{purpose}

Before that, we need two important results. These are presented in the next
two sections.

\subsection{Identifying superfluous monomial generators of the 
module quotients}
%\phantom{+}%%\bigskip

 Given a basis element $e_{k,i}=(I_1,\ldots
,I_k,I_{k+1})\in\mcb_k$, we want to identify superfluous elements
$m^k_{j,i}$ in $M_{i}(G_{k-1})$, where $1\leq j<i\leq r_k$ (recall
Notation~\ref{m-basis-elements} and Remark~\ref{uptosign}).

\begin{lemma}\label{superfluous}
Fix an integer $k$, with $1\leq k\leq n-2$. Suppose that $\mbfh_{k-1}$
holds. Fix a basis element $e_{k,i}=(I_1,\ldots ,I_k,I_{k+1})$ in
$\mcb_k$. Let $e_{k,j}=(J_1,\ldots ,J_k,J_{k+1})\in\mcb_k$, for some
$1\leq j<i\leq r_k$.
\begin{itemize}
\item[$(a)$] If for some $1\leq s\leq k-1$, $J_s\neq I_s$, then
  $m^k_{j,i}=0$. 
\end{itemize}
Suppose, on the contrary, that $e_{k,j}=(I_1,\ldots
,I_{k-1},J_{k},J_{k+1})$, so that $J_k\cup J_{k+1}=I_{k}\cup I_{k+1}$.
\begin{itemize}
\item[$(b)$] Then, 
\begin{eqnarray*}
m^k_{j,i} =(-1)^{k-1}x^{(J_k\cap I_k\to J_{k+1},I_{k+1})^+}x^{J_k\cap
  I_{k+1}\to J_{k+1}}.
\end{eqnarray*}
 In particular, if $J_k\supsetneq I_k$ (and so $J_{k+1}\subsetneq
 I_{k+1}$), then
\begin{eqnarray}\label{mji}
m^k_{j,i}=(-1)^{k-1}x^{J_{k}\cap I_{k+1}\to J_{k+1}}.
\end{eqnarray}
\item[$(c)$] If $J_k\not\supset I_k$, let $1\leq l<j$ be such that
  $e_{k,l}=(I_1,\ldots ,I_{k-1},J_{k}\cup I_k,J_{k+1}\cap
  I_{k+1})$. Then $m^k_{l,i}\mid m^k_{j,i}$, so $m^k_{j,i}$ is
  superfluous.
\end{itemize}
\end{lemma}
\begin{proof} 
Note that $(a)$ is vacuous if $k=1$. So, suppose momentarily, that
$k\geq 2$ and that $J_s\neq I_s$, for some $1\leq s\leq k-1$.  By the
hypothesis that $\mbfh_{k-1}$ holds, we have that
\begin{align}\label{lm-j-i}
 \lt(\dif_k(e_{k,j}))&=(-1)^{k-1}x^{J_k\to J_{k+1}}(J_1,\ldots
  ,J_{k-1},J_k\cup J_{k+1})
\intertext{and}
  \lt(\dif_k(e_{k,i}))&=(-1)^{k-1}x^{I_k\to I_{k+1}}(I_1,\ldots
  ,I_{k-1}, ,I_k\cup I_{k+1}), \notag
\end{align}
where $(J_1,\ldots ,J_k\cup J_{k+1})\neq (I_1,\ldots ,I_k\cup
I_{k+1})$. Then $(a)$ follows from the fact that the least common
multiple of two monomial terms is zero provided the two basis elements
involved are different (see Definition~\ref{firstsetofdefinitions},
$(2)$).

Suppose that $e_{k,j}=(I_1,\ldots ,I_{k-1},J_{k},J_{k+1})$, with
$J_k\cup J_{k+1}=I_{k}\cup I_{k+1}$. Note that this means that
$(J_k,J_{k+1})$ and $(I_k,I_{k+1})$ are each partitions of the same
set $J_k\cup J_{k+1}=I_{k}\cup I_{k+1}$. Now the two basis elements
involved in $\lm(\dif_{k}(e_{k,j}))$ and $\lm(\dif_k(e_{k,i}))$ are
the same (see \eqref{lm-j-i} above). Therefore, using
Definition~\ref{firstsetofdefinitions}, $(2)$, again:
\begin{align*}
m^k_{j,i} &= \frac{\Lcm(x^{J_k\to J_{k+1}},x^{I_k\to
      I_{k+1}})(I_1,\ldots ,I_{k-1},I_{k}\cup
    I_{k+1})}{(-1)^{k-1}x^{I_k\to I_{k+1}}(I_1,\ldots
    ,I_{k-1},I_{k}\cup I_{k+1})}
 \\ &= \frac{\Lcm(x^{J_k\to
      J_{k+1}},x^{I_k\to I_{k+1}})}{(-1)^{k-1}x^{I_k\to I_{k+1}}}.
\end{align*}
Now, using \eqref{disjoint}, we see that
\begin{eqnarray*}
&&x^{J_k\to J_{k+1}}=x^{J_k\cap I_{k}\to J_{k+1}}x^{J_{k}\cap I_{k+1}\to
  J_{k+1}}\mbox{ and }\\&&x^{I_k\to I_{k+1}}=x^{J_k\cap I_k\to
  I_{k+1}}x^{J_{k+1}\cap I_k\to I_{k+1}}.
\end{eqnarray*}
Thus
\begin{multline*}
\Lcm(x^{J_k\to J_{k+1}},x^{I_k\to I_{k+1}}) \\
=\Lcm(x^{J_k\cap I_{k}\to J_{k+1}}x^{J_{k}\cap I_{k+1}\to J_{k+1}},x^{J_k\cap
I_k\to I_{k+1}}x^{J_{k+1}\cap I_k\to I_{k+1}})\\
= \Lcm(x^{J_k\cap I_{k}\to J_{k+1}},x^{J_k\cap I_k\to I_{k+1}})
  x^{J_{k}\cap I_{k+1}\to J_{k+1}}x^{J_{k+1}\cap I_k\to I_{k+1}}.
\end{multline*}
Recalling \eqref{lcm-equality} in Notation~\ref{lcm-notation}, we
deduce:
\begin{eqnarray*}
\frac{\Lcm(x^{J_k\to J_{k+1}},x^{I_k\to I_{k+1}})}{(-1)^{k-1}x^{I_k\to
    I_{k+1}}}= (-1)^{k-1}x^{(J_k\cap I_k\to J_{k+1},I_{k+1})^+}
x^{J_k\cap I_{k+1}\to J_{k+1}}.
\end{eqnarray*}
In particular, if $J_k\supsetneq I_k$, then $J_{k+1}\subsetneq
I_{k+1}$ and $x^{(J_k\cap I_k\to J_{k+1},I_{k+1})^+}=1$. Therefore
\begin{eqnarray*}
m^k_{j,i}=(-1)^{k-1}x^{J_k\cap I_{k+1}\to J_{k+1}}.
\end{eqnarray*}
This proves $(b)$. Suppose now that $J_k\not\supset I_k$, so
$J_{k+1}\not\subset I_{k+1}$ and $J_{k+1}\cap I_{k+1}\subsetneq
J_{k+1}$. Let $1\leq l<j<i$ be such that $e_{k,l}=(I_1,\ldots
,I_{k-1},J_k\cup I_k,J_{k+1}\cap I_{k+1})$. Using $(b)$, we have that:
\begin{eqnarray*}
&& m^k_{l,i}=(-1)^{k-1}x^{J_k\cap I_{k+1}\to J_{k+1}\cap I_{k+1}}\mbox{
  and }\\ && m^k_{j,i}=(-1)^{k-1}x^{(J_k\cap I_k\to J_{k+1},I_{k+1})^+}
x^{J_k\cap I_{k+1}\to J_{k+1}}.
\end{eqnarray*}
In particular, $m^k_{l,i}$ divides $m^k_{j,i}$, which proves $(c)$.
\end{proof}

The following is an easy remark but will prove to be very useful in
organizing the basis elements of $\mcb_k$.

\begin{remark}\label{bki}
Fix an integer $k$, with $1\leq k\leq n-2$. Let $e_{k,i}=(I_1,\ldots
,I_k,I_{k+1})$ be a basis element of $\mcb_k$, where $1\leq i\leq
r_{k}$. Set
\begin{eqnarray*}
\mcb_{k,i}:=\{(I_1,\ldots ,I_{k-1},J_k,J_{k+1})\in\mcb_k\mid
J_k\supsetneq I_k\}.
\end{eqnarray*}
The following conditions hold.
\begin{itemize}
\item[$(a)$] The set $\mcb_{k,i}$ is empty if and only if
  $I_{k+1}=\{n\}$. (For instance, for the first basis element
\begin{eqnarray*}  
e_{k,1}=(\{k,\ldots ,n-1\},\{k-1\},\ldots ,\{2\},\{1\},\{n\})
\end{eqnarray*} 
of $\mcb_k$, considering the $\srle$, we have $\mcb_{k,1}=\emptyset$.)
\item[$(b)$] The cardinality of $\mcb_{k,i}$ is $\lvert
  \mcb_{k,i}\rvert=2^{\lvert I_{k+1}\rvert -1}-1$. (This formula holds
  trivially if $I_{k+1}=\{n\}$.)
\item[$(c)$] If $\mcb_{k,i}\neq \emptyset$ and $e_{k,j}=(I_1,\ldots
  ,I_{k-1},J_k,J_{k+1})$ in $\mcb_{k,i}$, then $1\leq j<i$. Moreover
  $\mcb_{k,j}\subsetneq\mcb_{k,i}$.
\item[$(d)$] If $\mcb_{k,i}\neq \emptyset$, let
  $\varrho_{k,i}:\mcb_{k,i}\to\mcb_{k+1}$ be defined by
\begin{eqnarray*}
\varrho_{k,i}(e_{k,j}):=(I_1,\ldots ,I_k,J_k\setminus I_k,J_{k+1}),
\end{eqnarray*}
while if $\mcb_{k,i}=\emptyset$, we set $\varrho_{k,i}(\mcb_{k,i}):=\emptyset$.
Then $\varrho_{k,i}$ is well-defined and injective.
\item[$(e)$] If $i\neq j$, then $\varrho_{k,i}(\mcb_{k,i})$ and
  $\varrho_{k,j}(\mcb_{k,j})$ are two disjoint sets.
\item[$(f)$] Then $\mcb_{k+1}=\bigcup_{i=2}^{r_k}\varrho_{k,i}(\mcb_{k,i})$.
In particular, $r_{k+1}=\sum_{i=2}^{r_k}\lvert\mcb_{k,i}\rvert$.
\end{itemize}
\end{remark}
\begin{proof} 
Suppose that $I_{k+1}=\{n\}$ and let $(I_1,\ldots
,I_{k-1},J_k,J_{k+1})\in\mcb_k$, with $J_k\supseteq I_k$. Then
\begin{eqnarray*}
J_k\cup J_{k+1}=[n]\setminus \cup_{s=1}^{k-1}I_s=I_k\cup I_{k+1}=I_k\cup
\{n\}. 
\end{eqnarray*}
Since $n\in J_{k+1}$, it follows that $J_k=I_k$ and so there
does not exist any element in $\mcb_{k,i}$. On the other hand, if
$I_{k+1}\supsetneq \{n\}$ and $m\in I_{k+1}$ with $m\neq n$, just take
$J_{k}=I_{k}\cup \{m\}$ and $J_{k+1}=I_{k+1}\setminus \{m\}$. Then
$(I_1,\ldots ,I_{k-1},J_k,J_{k+1})$ is in $\mcb_{k,i}$. This proves
$(a)$.

Suppose that $\mcb_{k,i}\neq \emptyset$, \ie $I_{k+1}\supsetneq
\{n\}$. So $I_k\subsetneq [n-1]\setminus \cup_{s=1}^{k-1}I_s$. Then
the elements $(I_1,\ldots ,I_{k-1},J_k,J_{k+1})$ of $\mcb_{k,i}$ are
in one-to-one correspondence with the non-empty subsets $J_k\setminus
I_k$ of the set $[n-1]\setminus \cup_{s=1}^kI_s=I_{k+1}\setminus
\{n\}$. This proves $(b)$.

Let $e_{k,j}=(I_1,\ldots ,I_{k-1},J_k,J_{k+1})\in\mcb_{k,i}$, that is,
$J_k\supsetneq I_k$. Then, according to the $\srle$, $e_{k,j}$
precedes $e_{k,i}$ and so $j<i$. Moreover, if $e_{k,h}=(I_1,\ldots
,I_{k-1},H_k,H_{k+1})$ is in $\mcb_{k,j}$, then $H_k\supsetneq
J_k\supsetneq I_k$, and so $e_{k,h}\in\mcb_{k,i}$. Furthermore,
$e_{k,j}\in\mcb_{k,i}$, but $e_{k,j}\not\in\mcb_{k,j}$. This proves
$(c)$.

Clearly, if $e_{k,j}=(I_1,\ldots ,I_{k-1},J_k,J_{k+1})\in\mcb_{k,i}$,
then $J_k\supsetneq I_k$ and $J_k\setminus I_k\neq \emptyset$. Thus
$\varrho_{k,i}(e_{k,j})$ is indeed in $\mcb_{k+1}$. Take
$e_{k,j}=(I_1,\ldots ,I_{k-1},J_k,J_{k+1})$ and $e_{k,h}=(I_1,\ldots
,I_{k-1},H_k,H_{k+1})$ in $\mcb_{k,i}$ such that
$\varrho_{k,i}(e_{k,j})=\varrho_{k,i}(e_{k,h})$. Then $J_{k+1}=H_{k+1}$
and, so $J_k=(J_{k}\cup J_{k+1})\setminus J_{k+1}= (H_k\cup
H_{k+1})\setminus H_{k+1}=H_k$, and $e_{k,j}=e_{k,h}$, which proves
$(d)$.

Next we prove $(e)$. Let $e_{k,i}=(I_1,\ldots ,I_k,I_{k+1})$ and
$e_{k,j}=(J_1,\ldots ,J_k,J_{k+1})$, with $\mcb_{k,i}\neq \emptyset$
and $\mcb_{k,j}\neq \emptyset$, respectively.  Suppose that
$e_{k,p}=(I_1,\ldots ,I_{k-1},P_k,P_{k+1})$ is in $\mcb_{k,i}$ and
$e_{k,q}=(J_1,\ldots ,J_{k-1},Q_k,Q_{k+1})$ is in
$\mcb_{k,j}$. Suppose also that
$\varrho_{k,i}(e_{k,p})=\varrho_{k,j}(e_{k,q})$. Then $I_s=J_s$, for
$s=1,\ldots ,k$. In particular, $I_{k+1}=[n]\setminus
\cup_{s=1}^kI_s=[n]\setminus\cup_{s=1}^kJ_s=J_{k+1}$. Thus
$e_{k,i}=e_{k,j}$. The rest follows by $(d)$.

Finally, we prove that $(f)$ holds. Consider $(H_1,\ldots
,H_{k+2})\in\mcb_{k+1}$. Let us find two basis elements
$e_{k,i}=(I_1,\ldots ,I_{k+1})\in\mcb_k$ and $e_{k,j}$ in
$\mcb_{k,i}$, so that $e_{k,j}=(I_1,\ldots ,I_{k-1},J_k,J_{k+1})$,
with $J_k\supsetneq I_k$, and such that
\begin{eqnarray*}
(H_1,\ldots ,H_{k+2})=\varrho_{k,i}(e_{k,j})=(I_1,\ldots
  ,I_k,J_k\setminus I_k,J_{k+1}).
\end{eqnarray*}
This forces $I_{s}=H_s$, for $s=1,\ldots ,k$, $J_{k}\setminus
I_k=H_{k+1}$ and $J_{k+1}=H_{k+2}$.  Therefore $I_{k+1}$ must be equal
to $[n]\setminus \cup_{s=1}^kI_s=[n]\setminus \cup_{s=1}^kH_s=H_{k+1}\cup
H_{k+2}$ and $J_k=I_k\cup H_{k+1}=H_k\cup H_{k+1}$.  Thus, both
\begin{eqnarray*}
&& e_{k,i}=(I_1,\ldots ,I_k,I_{k+1})=(H_1,\ldots ,H_k,H_{k+1}\cup
H_{k+2})\in\mcb_k\mbox{ and }\\ && e_{k,j}=(I_1,\ldots
,I_{k-1},J_k,J_{k+1})=(H_1,\ldots ,H_{k-1},H_{k}\cup
H_{k+1},H_{k+2})\in\mcb_{k,i},
\end{eqnarray*}
are uniquely determined by $(H_1,\ldots ,H_{k+2})\in\mcb_{k+1}$.
\end{proof}

Having in mind Notation~\ref{m-basis-elements},
Theorem~\ref{algorithm} and Remark~\ref{nonminimal}, then the
preceding Lemma~\ref{superfluous} and Remark~\ref{bki} can be
summarised as follows. Note that Lemma~\ref{superfluous} discards
elements $m^k_{j,i}$ which are known to be superfluous for
certain. However, it does not ensure that we have discarded all of
them. In fact, we will see that this is precisely characterised by the
digraph $\mbg$ being strongly complete (see
Corollary~\ref{min-resolution}).

\begin{proposition}\label{gbdesiredform}
Fix an integer $k$, with $1\leq k\leq n-2$. Suppose that $\mbfh_{k-1}$
holds. For each $i=2,\ldots ,r_k$, then
\begin{eqnarray*}
\{ m^k_{j,i}\mid e_{k,j}\in\mcb_{k,i}\}
\end{eqnarray*}
is a generating set for the module quotient $M_i(G_{k-1})$. Let
$\tau_k(e_{k,i},e_{k,j})$ stand for a $\tau$-syzygy associated to the
$S$-vector of $f_{k-1,i}=\dif_k(e_{k,i})$ and
$f_{k-1,j}=\dif_k(e_{k,j})$, where $e_{k,j}\in\mcb_{k,i}$. Then there
is a Gr\"obner basis of $\ker(\dif_k)\subset\mcc_k$ (w.r.t. the
monomial ordering on $\mcc_k$ induced by the monomial ordering on
$\mcc_{k-1}$ and the Gr\"obner basis $G_{k-1}$) of the form
\begin{eqnarray}\label{gbdf}
\bigcup_{i=2}^{r_k}\{\tau_{k}(e_{k,i},e_{k,j})\mid
e_{k,j}\in\mcb_{k,i}\}.
\end{eqnarray}
The cardinality of this Gr\"obner basis is
$\sum_{i=2}^{r_k}\lvert\mcb_{k,i}\rvert$, which coincides with
$r_{k+1}$, the rank of $\mcc_{k+1}$.
\end{proposition}
\begin{proof}
Indeed, Lemma~\ref{superfluous} ensures that in order to find a
generating set of $M_i(G_{k-1})$ we only need to consider the
$m^{k}_{j,i}$ with $e_{k,j}\in\mcb_{k,i}$. By Theorem~\ref{algorithm}
and Remark~\ref{nonminimal}, we deduce that indeed there is a Gr\"obner
basis of $\ker(\dif_k)$ of the form \eqref{gbdf}.

By Discussion~\ref{disc}, we know that if
$\tau_{k}(e_{k,i},e_{k,j})=\tau_k(e_{k,p},e_{k,q})$, then
$e_{k,i}=e_{k,p}$ and $e_{k,j}=e_{k,q}$. In particular, the union in
\eqref{gbdf} is of disjoint sets and
\begin{eqnarray*}
\lvert\{\tau_k(e_{k,i},e_{k,j})\mid
e_{k,j}\in\mcb_{k,i}\}\rvert=\lvert\mcb_{k,i}\rvert. 
\end{eqnarray*}
Thus, the cardinality of this Gr\"obner basis is
$\sum_{i=2}^{r_k}\lvert\mcb_{k,i}\rvert$. By Remark~\ref{bki}, $(f)$,
$\sum_{i=2}^{r_k}\lvert\mcb_{k,i}\rvert=r_{k+1}$.
\end{proof}

\begin{remark}\label{min-m}
With the notations above, if $\mbg$ is strongly complete or,
equivalently, if $L$ is a $\pcb$ matrix, then
\begin{eqnarray*}
\vert\mcb_{k,i}\rvert=
\lvert \{ m^k_{j,i}\mid e_{k,j}\in\mcb_{k,i}\}\rvert.
\end{eqnarray*}
For in the strongly complete case, the map $\mcb_{k,i}\to\mcc_{k-1}$
defined by $e_{k,j}\mapsto m^{k}_{j,i}$ is injective. Indeed, if
$e_{k,j}\in\mcb_{k,i}$, then $m^k_{j,i}=(-1)^{k-1}x^{J_{k}\cap
  I_{k+1}\to J_{k+1}}$ (see \eqref{mji} in Lemma~\ref{superfluous}),
where $e_{k,i}=(I_1,\ldots ,I_{k+1})$ and $e_{k,j}=(I_1,\ldots,
I_{k-1},J_k,J_{k+1})$, with $J_k\supsetneq I_k$. Since $a_{i,j}>0$,
for all $i,j$, it follows that $m^k_{j,i}$ determines $J_k\cap
I_{k+1}$. Hence $J_k=(J_k\cap I_k)\cup (J_k\cap I_{k+1})=I_k\cup
(J_k\cap I_{k+1})$, and so $J_{k+1}$ and hence $e_{k,j}$ are uniquely
determined by $m^k_{j,i}$.
\end{remark}

\subsection{The images  of basis elements 
are syzygies associated to {\em S}-vectors} %\phantom{+}%%\bigskip

 Let us show that the differential $\dif_{k+1}$ of the $\cyc$
complex maps each basis element of $\mcb_{k+1}$ to a $\tau$-syzygy
associated to an $S$-vector, up to a $\pm$ sign (see
Discussion~\ref{disc} ).

\begin{proposition}\label{tau}
Fix an integer $k$, with $1\leq k\leq n-2$. Suppose that $\mbfh_{k-1}$
holds. Fix a basis element $(I_1,\ldots ,I_{k+2})$ in $\mcb_{k+1}$.
Then
\begin{eqnarray*}
\dif_{k+1}(I_1,\ldots,I_{k+2})=-\tau_{k}(e_{k,i},e_{k,j}),
\end{eqnarray*}
where $e_{k,i}=(I_1,\ldots ,I_{k+1}\cup I_{k+2})$ and
$e_{k,j}=(I_1,\ldots,I_{k}\cup I_{k+1},I_{k+2})$. That is,
$-\dif_{k+1}(I_1,\ldots,I_{k+2})$ is a $\tau$-syzygy associated to the
$S$-vector of $f_{k-1,i}=\dif_{k}(e_{k,i})$ and
$f_{k-1,j}=\dif_{k}(e_{k,j})$.
\end{proposition}
\begin{proof}
Note that, according to the $\srle$, $e_{k,j}$ does indeed precede
$e_{k,i}$ (see Remark~\ref{srle}) and it follows from the definition
that $e_{k,j}\in\mcb_{k,i}$ (see Remark~\ref{bki}). By
Definition~\ref{cyccomplex}, we have:
\begin{align*}
\dif_{k+1}(I_1,\ldots,I_{k+2})=& (-1)^{k}x^{I_{k+1}\to
    I_{k+2}}(I_1,\ldots ,I_{k+1}\cup I_{k+2})\\
&+ (-1)^{k-1}x^{I_{k}\to I_{k+1}}(I_1,\ldots ,I_{k}\cup
  I_{k+1},I_{k+2})\\
&+\sum_{s=1}^{k-1}(-1)^{s-1}x^{I_s\to
    I_{s+1}}(I_1,\ldots ,I_s\cup I_{s+1},\ldots
  ,I_{k+2})\\
&-x^{I_{k+2}\to I_1}(I_2,\ldots,I_{k+1},I_1\cup I_{k+2}).
\end{align*}
Set $e_{k,j_s}:=(I_1,\ldots ,I_s\cup I_{s+1},\ldots ,I_{k+2})$, for
some $j_s$, with $1\leq j_1<\cdots< j_s<\cdots<j_{k-1}<j<i$.
Moreover, set $e_{k,l}:=(I_2,\ldots,I_{k+1},I_1\cup I_{k+2})$, for
some $l\in \{1,\ldots ,r_{k}\}$. Then
\begin{multline}\label{difbasisele}
 \dif_{k+1}(I_1,\ldots,I_{k+2})= (-1)^{k}x^{I_{k+1}\to
  I_{k+2}}e_{k,i}\\ 
+(-1)^{k-1}x^{I_{k}\to   I_{k+1}}e_{k,j}+ \sum_{s=1}^{k-1}(-1)^{s-1}x^{I_s\to
  I_{s+1}}e_{k,j_s}-x^{I_{k+2}\to I_1}e_{k,l}.
\end{multline}
Since $\dif_{k}\circ\dif_{k+1}=0$, and $\dif_{k}(e_{k,u})=f_{k-1,u}$,
for each $u$, we obtain:
\begin{multline}\label{s-vector}
 (-1)^{k-1}x^{I_{k+1}\to I_{k+2}}f_{k-1,i}-(-1)^{k-1}x^{I_{k}\to
    I_{k+1}}f_{k-1,j}\\
= \sum_{s=1}^{k-1}(-1)^{s-1}x^{I_s\to
    I_{s+1}}f_{k-1,j_s}-x^{I_{k+2}\to I_1}f_{k-1,l}.
\end{multline}
Let us see that the left hand side of \eqref{s-vector} coincides
with the $S$-vector of $f_{k-1,i}$ and $f_{k-1,j}$. 

Since $\mbfh_{k-1}$ holds, we have
\begin{align*}
\lt(f_{k-1,j})&=\lt(\dif_{k}(e_{k,j}))=\lt(\dif_{k}(I_1,\ldots,I_{k}\cup
I_{k+1},I_{k+2}))\\
&= (-1)^{k-1}x^{I_{k}\cup I_{k+1}\to I_{k+2}}(I_1,\ldots
,I_{k}\cup I_{k+1}\cup I_{k+2}).
\end{align*}
Similarly,
\begin{align*}
\lt(f_{k-1,i})&=\lt(\dif_{k}(e_{k,i}))=\lt(\dif_{k}(I_1,\ldots,
I_{k+1}\cup I_{k+2}))\\
&= (-1)^{k-1}x^{I_{k}\to I_{k+1}\cup I_{k+2}}(I_1,\ldots
,I_{k}\cup I_{k+1}\cup I_{k+2}).
\end{align*}
Then 
\begin{multline*}
\Lcm(\lm(f_{k-1,j}),\lm(f_{k-1,i}))\\
= \Lcm(x^{I_{k}\cup I_{k+1}\to
  I_{k+2}},x^{I_{k}\to I_{k+1}\cup I_{k+2}})(I_1,\ldots ,I_{k}\cup
I_{k+1}\cup I_{k+2}).
\end{multline*}
Using \eqref{disjoint} in Notation~\ref{lcm-notation},
\begin{eqnarray*}
x^{I_{k}\cup I_{k+1}\to I_{k+2}}=x^{I_{k}\to I_{k+2}}x^{I_{k+1}\to
  I_{k+2}}\mbox{ and } x^{I_{k}\to I_{k+1}\cup I_{k+2}}=x^{I_{k}\to
  I_{k+1}}x^{I_{k}\to I_{k+2}}.
\end{eqnarray*}
Therefore,
\begin{eqnarray*}
\Lcm(x^{I_{k}\cup I_{k+1}\to I_{k+2}},x^{I_{k}\to I_{k+1}\cup
  I_{k+2}})=x^{I_{k}\to I_{k+2}}x^{I_{k+1}\to I_{k+2}}x^{I_{k}\to I_{k+1}}.
\end{eqnarray*}
Hence
\begin{eqnarray*}
m^{k}_{j,i}= \frac{\Lcm(x^{I_{k}\cup I_{k+1}\to
    I_{k+2}},x^{I_{k}\to I_{k+1}\cup
    I_{k+2}})}{(-1)^{k-1}x^{I_{k}\to I_{k+1}\cup
    I_{k+2}}}=(-1)^{k-1}x^{I_{k+1}\to I_{k+2}},
\end{eqnarray*}
and, similarly, $m^{k}_{i,j}=(-1)^{k-1}x^{I_{k}\to
  I_{k+1}}$. Therefore, the $S$-vector of $f_{k-1,i}$ and $f_{k-1,j}$ is
equal to
\begin{align*}
 S(f_{k-1,i},f_{k-1,j})&=m^{k}_{j,i}f_{k-1,i}-m^{k}_{i,j}f_{k-1,j}\\
&= (-1)^{k-1}x^{I_{k+1}\to
  I_{k+2}}f_{k-1,i}-(-1)^{k-1}x^{I_{k}\to I_{k+1}}f_{k-1,j},
\end{align*}
which is the left hand side of \eqref{s-vector}.

Now, let us prove that the right hand side of \eqref{s-vector} is a
standard expression for the $S$-vector $S(f_{k-1,i},f_{k-1,j})$
w.r.t. the Gr\"obner basis $G_{k-1}=\{f_{k-1,1},\ldots
,f_{k-1,r_{k}}\}$, with remainder zero. Note that $G_{k-1}$ is indeed
a Gr\"obner basis of $\ker(\dif_{k-1})$ due to hypothesis
$\mbfh_{k-1}$.  First, observe that the only coincident terms among
the summands of
\begin{eqnarray*}
m^{k}_{j,i}f_{k-1,i}=(-1)^{k-1}x^{I_{k+1}\to I_{k+2}}f_{k-1,i}\;\mbox{ and }\;
m^{k}_{i,j}f_{k-1,j}=(-1)^{k-1}x^{I_{k}\to I_{k+1}}f_{k-1,j}
\end{eqnarray*}
are precisely the leading terms. Indeed: 
\begin{align}\label{fi}
m^{k}_{j,i}f_{k-1,i}&=(-1)^{k-1}x^{I_{k+1}\to
  I_{k+2}}f_{k-1,i}=(-1)^{k-1}x^{I_{k+1}\to
  I_{k+2}}\dif_{k}(e_{k,i})\notag \\ 
&= (-1)^{k-1}x^{I_{k+1}\to
  I_{k+2}}\dif_{k}(I_1,\ldots ,I_{k+1}\cup
I_{k+2}) \notag \\
&= (-1)^{k-1}x^{I_{k+1}\to
  I_{k+2}}\sum_{s=1}^{k-1}(-1)^{s-1}x^{I_s\to I_{s+1}}(I_1,\ldots\\
&\phantom{(-1)^{k-1}x^{I_{k+1}\to
  I_{k+2}}\lim_{s=1}(-1)^{s-1}x^{I_s\to I_{s+1}}}\dots,I_s\cup I_{s+1},\ldots ,I_{k+1}\cup I_{k+2})\notag
\\
&+ \underline{(-1)^{k-1}x^{I_{k+1}\to I_{k+2}}(-1)^{k-1}x^{I_{k}\to
    I_{k+1}\cup I_{k+2}}(I_1,\ldots ,I_{k}\cup I_{k+1}\cup
  I_{k+2})}\notag \\ 
& -(-1)^{k-1}x^{I_{k+1}\to I_{k+2}}x^{I_{k+1}\cup I_{k+2}\to
  I_1}(I_2,\ldots,I_1\cup I_{k+1}\cup I_{k+2})\notag
\end{align}
whereas 
\begin{align}\label{fj}
 m^{k}_{i,j}f_{k-1,j}&=(-1)^{k-1}x^{I_{k}\to I_{k+1}}f_{k-1,j}=
  (-1)^{k-1}x^{I_{k}\to I_{k+1}}\dif_{k}(e_{k,j}) \notag \\ 
&=  (-1)^{k-1}x^{I_{k}\to I_{k+1}}\dif_{k}(I_1,\ldots ,I_{k}\cup
  I_{k+1},I_{k+2}) \notag \\
&= (-1)^{k-1}x^{I_{k}\to
    I_{k+1}}\sum_{s=1}^{k-2}(-1)^{s-1}x^{I_s\to I_{s+1}}(I_1,\ldots
  \\
	&\phantom{(-1)^{k-1}x^{I_{k}\to
    I_{k+1}}\lim_{s=1}(-1)^{s-1}x^{I_s\to I_{s+1}}}\ldots, I_s\cup I_{s+1}, \dots I_{k}\cup I_{k+1}, I_{k+2})\notag\\ 
&+  (-1)^{k-1}x^{I_{k}\to I_{k+1}}(-1)^{k-2}x^{I_{k-1}\to I_{k}\cup
    I_{k+1}} (I_1,\ldots , I_{k-1}\cup I_{k}\cup I_{k+1}, I_{k+2})\notag
 \\
&+ \underline{(-1)^{k-1}x^{I_{k}\to
      I_{k+1}}(-1)^{k-1}x^{I_{k}\cup I_{k+1}\to I_{k+2}}(I_1,\ldots
    ,I_{k}\cup I_{k+1}\cup I_{k+2})} \notag \\
  &-(-1)^{k-1}x^{I_{k}\to I_{k+1}}x^{I_{k+2}\to
    I_1}(I_2,\ldots,I_{k}\cup I_{k+1},I_1\cup I_{k+2})\notag
\end{align}
Excluding the (underlined) coincident leading terms (note that
hypothesis $\mbfh_{k-1}$ holds), the basis elements involved in
\eqref{fi} end either in $I_{k+1}\cup I_{k+2}$ or else in $I_1\cup I_{k+1}\cup
I_{k+2}$, while the basis elements involved in \eqref{fj} end either
in $I_{k+2}$ or else in $I_1\cup I_{k+2}$. Therefore, the $S$-vector
of $f_{k-1,i}$ and $f_{k-1,j}$ cancels the leading terms of these
expressions, but no other summand of $m^{k}_{j,i}f_{k-1,i}$ is
cancelled by any other summand of $m^{k}_{i,j}f_{k-1,j}$.  In
particular,
\begin{eqnarray}\label{geq1}
\lm(S(f_{k-1,i},f_{k-1,j}))\geq
\lm(m^{k}_{j,i}f_{k-1,i}-\lt(m^{k}_{j,i}f_{k-1,i})).
\end{eqnarray}

Now, we need to identify the leading monomials of the summands on the
right hand side of \eqref{s-vector}. For $1\leq s\leq k-1$, since the
hypothesis $\mbfh_{k-1}$ holds,
\begin{align}\label{lmRHSs}
\lm(x^{I_s\to I_{s+1}}f_{k-1,j_s})&=x^{I_s\to
    I_{s+1}}\lm(\dif_{k}(e_{k,j_s}))\\
&= x^{I_s\to  I_{s+1}}\lm(\dif_{k}(I_1,\ldots ,I_s\cup I_{s+1},\ldots
  ,I_{k+2}))\notag \\ 
&= x^{I_s\to I_{s+1}}x^{I_{k+1}\to I_{k+2}}
  (I_1,\ldots ,I_s\cup I_{s+1},\ldots ,I_{k+1}\cup I_{k+2}). \notag
\end{align}
Similarly, the leading monomial of the last summand on the right hand
side of \eqref{s-vector} is:
\begin{align}\label{lmRHSl}
\lm(x^{I_{k+2}\to I_{1}}f_{k-1,l})&=x^{I_{k+2}\to
    I_{1}}\lm(\dif_{k}(e_{k,l}))\\
&= x^{I_{k+2}\to     I_{l}}\lm(\dif_{k}(I_2,\ldots ,I_{k+1}, I_{1}\cup I_{k+2}))
  \notag \\
&= x^{I_{k+2}\to I_{1}}x^{I_{k+1}\to I_1\cup I_{k+2}}
  (I_2,\ldots ,I_1\cup I_{k+1}\cup I_{k+2}). \notag
\end{align}
All the monomials terms in \eqref{lmRHSs} and in \eqref{lmRHSl} appear
in \eqref{fi} and are different from the underlined leading term
$\lt(m^{k}_{j,i}f_{k-1,i})$. (By \eqref{disjoint}, $x^{I_{k+2}\to
  I_1}x^{I_{k+1}\to I_1\cup I_{k+2}}=x^{I_{k+1}\to
  I_{k+2}}x^{I_{k+1}\cup I_{k+2}\to I_1}$.)  Thus
\begin{equation}\label{geq2}
\lm(m^{k}_{j,i}f_{k-1,i}-\lt(m^{k}_{j,i}f_{k-1,i}))\geq\lm(x^{I_s\to I_{s+1}}f_{k-1,j_s}),\lm(x^{I_{k+2}\to I_{1}}f_{k-1,l}).
\end{equation}
Concatenating \eqref{geq1} and \eqref{geq2}, we deduce that the right
hand side of \eqref{s-vector} is indeed a standard expression for
$S(f_{k-1,i},f_{k-1,j})$ with remainder zero and w.r.t. the Gr\"obner
basis $G_{k-1}$.

Hence, the corresponding $\tau$-syzygy $\tau_{k}(e_{k,i},e_{k,j})$
associated to the $S$-vector of the elements
$f_{k-1,i}=\dif_{k}(e_{k,i})$ and $f_{k-1,j}=\dif_{k}(e_{k,j})$ is
\begin{eqnarray}\label{tau-vector}
&& \tau_{k}(e_{k,i},e_{k,j})=(-1)^{k-1}x^{I_{k+1}\to
  I_{k+2}}e_{k,i}\\ && - (-1)^{k-1}x^{I_{k}\to
  I_{k+1}}e_{k,j} -\sum_{s=1}^{k-1}(-1)^{s-1}x^{I_s\to
  I_{s+1}}e_{k,j_s}+x^{I_{k+2}\to I_1}e_{k,l}. \notag
\end{eqnarray}
Comparing \eqref{tau-vector} with \eqref{difbasisele}, we get
$\dif_{k+1}(I_1,\ldots,I_{k+2})=-\tau_{k}(e_{k,i},e_{k,j})$, the
desired equality. 
\end{proof}

As an immediate consequence of Proposition~\ref{tau}, we get the
following result. The second part can be seen as a natural
generalisation of Lemma~\ref{fc=fd}.

\begin{corollary}\label{ltbij}
Fix an integer $k$, with $1\leq k\leq n-2$. Suppose that $\mbfh_{k-1}$
holds. Then 
\begin{itemize}
\item[$(a)$]
  $\lt(\dif_{k+1}(I_1,\ldots,I_{k+1},I_{k+2}))=(-1)^{k}x^{I_{k+1}\to
  I_{k+2}}(I_1,\ldots, I_{k},I_{k+1}\cup I_{k+2})$.
\item[$(b)$] Moreover, the restriction of the map
  $\dif_{k+1}:\mcc_{k+1}\to\mcc_{k}$ over the basis set $\mcb_{k+1}$
  is injective. In particular, since $G_k:=\dif_{k+1}(\mcb_{k+1})$,
  then 
\begin{eqnarray*}
\lvert
G_k\rvert=\lvert\dif_{k+1}(\mcb_{k+1})\rvert=
\lvert\mcb_{k+1}\rvert=r_{k+1}.
\end{eqnarray*}
\end{itemize}
\end{corollary}
\begin{proof}
By Proposition~\ref{tau}, and with the notations used there, and by
\eqref{inducedlm} in Discussion~\ref{disc},
\begin{align*}
\lt(\dif_{k+1}(I_1,\ldots,I_{k+2}))&=-\lt(\tau_{k}(e_{k,i},e_{k,j}))
=-m^{k}_{j,i}e_{k,i}
\\ &= (-1)(-1)^{k-1}x^{I_{k+1}\to I_{k+2}}(I_1,\ldots,
I_{k+1}\cup I_{k+2})\\
&= (-1)^{k}x^{I_{k+1}\to I_{k+2}}(I_1,\ldots,
I_{k+1}\cup I_{k+2}),
\end{align*}
which proves $(a)$.

Suppose that $\dif_{k+1}(I_1,\ldots ,I_{k+2})=\dif_{k+1}(J_1,\ldots
,J_{k+2})$. Hence their leading terms coincide. Using the first part
$(a)$, it follows that
\[
 (-1)^{k}x^{I_{k+1}\to I_{k+2}}(I_1,\ldots,I_{k},I_{k+1}\cup
  I_{k+2})\!=\! (-1)^{k}x^{J_{k+1}\to
    J_{k+2}}(J_1,\ldots,J_{k},J_{k+1}\cup J_{k+2}).
\]
Therefore, $I_s=J_s$, for every $1\leq s\leq k$, and moreover,
$I_{k+1}\cup I_{k+2}=J_{k+1}\cup J_{k+2}$. Consider the summand
\begin{eqnarray*}
(-1)^{k-1}x^{I_{k}\to I_{k+1}}(I_1,\ldots ,I_{k-1},I_{k}\cup
  I_{k+1},I_{k+2})
\end{eqnarray*} 
of $\dif_{k+1}(I_1,\ldots ,I_{k+2})$, preceding the leading term and
different from 
\begin{eqnarray*}
-x^{I_{k+2}\to I_1}(I_2,\ldots ,I_{k+1},I_1\cup I_{k+2})
\end{eqnarray*}
(see \eqref{differential} in Definition~\ref{cyccomplex} and
Remark~\ref{dif-and-enum}). On the one hand, this element coincides
with
\begin{eqnarray*}
(-1)^{k-1}x^{I_{k}\to I_{k+1}}(J_1,\ldots ,J_{k-1},I_{k}\cup
  I_{k+1},I_{k+2}). 
\end{eqnarray*}
On the other hand, this element must coincide with a summand of
$\dif_{k+1}(J_1,\ldots ,J_{k+2})$. It follows that either $I_{k}\cup
I_{k+1}=J_{k}\cup J_{k+1}$ and $I_{k+2}=J_{k+2}$, or else, $I_{k}\cup
I_{k+1}=J_{k}$ and $I_{k+2}=J_{k+1}\cup J_{k+2}$. However, the second
case cannot occur, because, since $I_{k}=J_{k}$, it would follow that
$I_{k+1}=\emptyset$, a contradiction. Thus the first case holds, which
clearly implies $I_{k+1}=J_{k+1}$ and $I_{k+2}=J_{k+2}$.
\end{proof}

\subsection{Proof of the main theorem}
%\phantom{+}%%\bigskip

 Now, we have all the ingredients to prove
Theorem~\ref{main}.

